
\documentclass[a4paper,10pt]{report}

%\usepackage[landscape]{geometry}
\usepackage[cm-default]{fontspec}
\setromanfont{FreeSerif}
\setsansfont{FreeSans}
\setmonofont{FreeMono}
\usepackage[utf8x]{inputenc}
\usepackage{fontspec}
\usepackage{xunicode}
\usepackage{xltxtra}
\usepackage{xgreek}
\usepackage{amsmath}
\usepackage{unicode-math}
\usepackage{ulem}
\usepackage{color}
\usepackage{verbatim}
\usepackage{nopageno}
\usepackage{graphicx}
\usepackage{textpos}
\setlength{\TPHorizModule}{1cm}
\setlength{\TPVertModule}{1cm}
%\usepackage[colorgrid,texcoord]{eso-pic}
\usepackage[outline]{contour}
\usepackage{wrapfig}
\usepackage{url}
\usepackage{color}
\usepackage{tikz}
\usepackage{fancybox,fancyhdr}
\usepackage{subfigure}
\usepackage{pstricks}
\usepackage{epsfig}
\usepackage{multicol}
\usepackage{listings}
\usepackage{enumerate}
\usepackage{hyperref}
\hypersetup{
  bookmarks=true,
  bookmarksopen=true,
  pdfborder=false,
  pdfpagemode=UseNone,
  raiselinks=true,
  pdfhighlight={/P},
  colorlinks,
  citecolor=black,
  filecolor=black,
  linkcolor=black,
  urlcolor=black
}


\usepackage{fontspec}
\usepackage{xunicode}
\usepackage{xltxtra}


% Margins
%\setlength{\textwidth}{24cm}

%\setlength{\voffset}{0in}
%\setlength{\textheight}{6.8in}
% Colors
\definecolor{cbrown}{rgb}{0.49,0.24,0.07}
\definecolor{cmpez}{rgb}{0.92,0.88,0.79}
\definecolor{cmpez2}{rgb}{0.62,0.33,0.34}
\definecolor{cwhite}{rgb}{1,1,1}
\definecolor{cblack}{rgb}{0,0,0}
\definecolor{cred}{rgb}{0.9,0.15,0.15}
\definecolor{clblue}{rgb}{0.57,0.85,0.97}
\definecolor{corange}{rgb}{0.86,0.49,0.18}
\definecolor{clorange}{rgb}{0.95,0.84,0.61}
\definecolor{cyellow}{rgb}{1,0.95,0.16}
\definecolor{cgreen}{rgb}{0.69,0.8,0.31}
\definecolor{clbrown}{rgb}{0.78,0.67,0.43}
\definecolor{cmagenta}{rgb}{0.79,0.05,0.54}
\definecolor{cgray}{rgb}{0.76,0.73,0.63}
\definecolor{bluesite}{rgb}{0.05,0.43,0.67}
\definecolor{redsite}{rgb}{0.90,0.15,0.15}
\definecolor{backsite}{rgb}{0.07,0.07,0.07}
% ----------- from Costas -----------------------------

%---------------   Main Fonts    -----------------%   

% \setmainfont[Mapping=tex-text]{Calibri}
% \setmainfont[Mapping=tex-text]{Times New Roman}
%\setmainfont[Mapping=tex-text]{Droid Serif}
%\setmainfont[Mapping=tex-text]{Cambria}
%\setmainfont[Mapping=tex-text]{cm-unicode}
% \setmainfont[Mapping=tex-text]{Gentium}
% \setmainfont[Mapping=tex-text]{GFS Didot}
% \setmainfont[Mapping=tex-text]{Comic Sans MS}
 \setmainfont[Mapping=tex-text]{Ubuntu}
% \setmainfont[Mapping=tex-text]{Myriad Pro}
% \setmainfont[Mapping=tex-text]{CMU Concrete}
%\setmainfont[Mapping=tex-text]{DejaVu Sans}
%\setmainfont[Mapping=tex-text]{KerkisSans}
%\setmainfont[Mapping=tex-text]{KerkisCaligraphic}
% \setmainfont[Mapping=tex-text]{Segoe Print}
% \setmainfont[Mapping=tex-text]{Gabriola}

%------------------------------------------------%

% ------------ Mathematics Fonts ----------------%
\setmathfont{Asana-Math.ttf}
% -----------------------------------------------%
% Misc
\author{Κασωτάκη Ε. - Σμαραγδάκης Κ.}
\title{www.math24.gr}
% ---------------- Formation --------------------%


\setlength\topmargin{-0.7cm}
\addtolength\topmargin{-\headheight}
\addtolength\topmargin{-\headsep}
\setlength\textheight{26cm}
\setlength\oddsidemargin{-0.54cm}
\setlength\evensidemargin{-0.54cm}
%\setlength\marginparwidth{1.5in}
\setlength\textwidth{17cm}

\RequirePackage[avantgarde]{quotchap}
\renewcommand\chapterheadstartvskip{\vspace*{0\baselineskip}}
\RequirePackage[calcwidth]{titlesec}
\titleformat{\section}[hang]{\bfseries}
{\Large\thesection}{12pt}{\Large}[{\titlerule[0.9pt]}]
%--------------------------------------------------%
\usepackage{draftwatermark}
\SetWatermarkText{www.math24.gr}
\SetWatermarkLightness{0.9}
\SetWatermarkScale{0.6}
%--------------------------------------------------%


\usepackage{xcolor}
\usepackage{amsthm}
\usepackage{framed}
\usepackage{parskip}

\colorlet{shadecolor}{bluesite!20}

\newtheorem*{orismos}{Ορισμός}

\newenvironment{notation}
  {\begin{shaded}\begin{theorem}}
  {\end{theorem}\end{shaded}}




% Document begins
\begin{document}
\pagestyle{fancy}
\fancyhead{}
\fancyfoot{}
\renewcommand{\headrulewidth}{0pt}
\renewcommand{\footrulewidth}{0pt}


\fancyhead[LO,LE]{
 %\textblockcolor{backsite}
 %\begin{textblock}{5}(-2,-0.55)
  %\rule{0cm}{1cm}
 %\end{textblock}
 \textblockcolor{bluesite}
 \begin{textblock}{5}(-1.5,-0.55)
  \rule{0cm}{1cm}
 \end{textblock}
 %\textblockcolor{bluesite}
 %\begin{textblock}{14}(5.5,-0.55)
 % \rule{0cm}{1cm}
 %\end{textblock}
 \begin{textblock}{0}(-1,-0.25)
 \color{cwhite} \begin{Large}www.math24.gr\end{Large}
 \end{textblock}
\begin{textblock}{0}(16,-0.4)
 \color{cwhite} \includegraphics[height=1.5cm]{math24_logo.png}
 \end{textblock}
\textblockcolor{white}
\begin{textblock}{17}(-1,28)
 \color{backsite} \begin{small}Copyright \textcopyright 2012, Κασωτάκη Ε.(ikasotaki@gmail.com) - Σμαραγδάκης Κ.(kesmarag@gmail.com)\end{small}
 \end{textblock}
}
 
\begin{shaded}
\begin{center}
\huge \textbf{Φυλλάδιο Ασκήσεων}\\
%Πρόσθεση ρητών αριθμών
\end{center} 
\textbf{Μαθηματικά Α' Γυμνασίου} \hfill \textbf{Ημερομηνία Παράδοσης : \hspace{2em} }
\subsection*{Ονοματεπώνυμο :\hfill  \hspace{5em}}
\end{shaded}
\vspace{2em}
\begin{itemize}
 \item Φυσικοί αριθμοί
 \item Πρόσθεση φυσικών αριθμών
 \item Ιδιότητες της πρόσθεσης
 \item Αφαίρεση φυσικών αριθμών
\end{itemize}
\section*{Θεωρία - Φυσικοί αριθμοί\hfill \small{}}
\begin{itemize}
 \item \textbf{Φυσικοί αριθμοί} είναι οι αριθμοί $0,1,2,3,4,\cdots ,11,12,\cdots , 997, \cdots$
 \item Κάθε φυσικός αριθμός έχει έναν επόμενο και έναν προηγούμενο αριθμό (εκτός από το $0$ που έχει 
       μόνο επόμενο, το $1$)
\end{itemize}

\section*{Θεωρία - Πρόσθεση Φυσικών Αριθμών \hfill \small{}}
\textbf{Πρόσθεση} είναι η πράξη με την οποία από δύο φυσικούς αριθμούς $α$ και $β$ (τους προσθετέους), 
βρίσκουμε ένα τρίτο φυσικό αριθμό $γ$ που είναι το άθροισμά τους. Δηλαδή $α+β=γ$\\
\textbf{π.χ} $3+4=7$ (προσθετέοι: $3$ και $4$, άθροισμα: $7$)\\
\textbf{π.χ} $8+1=9$ (προσθετέοι: $8$ και $1$, άθροισμα: $9$)

\section*{Άσκηση 1  \hfill \small{20 μονάδες}}
Να υπολογίσετε τα παρακάτω αθροίσματα: 
\begin{enumerate} [1)]
 \item $13+11$
 \item $20+40$
 \item $4+14$
 \item $19+37$
 \item $16+751$
\end{enumerate}

\section*{Θεωρία - Ιδιότητες της Πρόσθεσης \hfill \small{}}
\textbf{Ιδιότητες της πρόσθεσης:}
\begin{itemize}
 \item \textbf{Πρόσθεση με το $0$}: το $0$ όταν προστεθεί σε ένα φυσικό αριθμό δεν τον μεταβάλλει,\\
        δηλαδή $0+α=α$ και $α+0=α$\\
        \textbf{π.χ} $19+0=19$\\
        \textbf{π.χ} $0+761=761$
 \item \textbf{Αντιμεταθετική ιδιότητα: } Μπορούμε να αλλάξουμε τη σειρά των δύο προσθετέων ενός αθροίσματος,\\
        δηλαδή $α+β=β+α$\\
       \textbf{π.χ} $3+4=4+3$ (γιατί $3+4=7$ και $4+3=7$)\\
       \textbf{π.χ} $11+25=25+11$ (γιατί $11+25=36$ και $25+11=36$)
 \item \textbf{Προσεταιριστική ιδιότητα:} Μπορούμε να αντικαθιστούμε προσθετέους με το άθροισμα τους ή να αναλύουμε 
       ένα προσθετέο σε άθροισμα,\\
       δηλαδή $α+β+γ=α+(β+γ)=(α+β)+γ$\\
       \textbf{π.χ} $11+(9+8)=(11+9)+8$ (γιατί $11+(9+8)=11+17=28$ και $(11+9)+8=20+8=28$)\\
       \textbf{π.χ} $1+(2+3)=(1+2)+3$(γιατί $1+(2+3)=1+5=6$ και $(1+2)+3=3+3=6$)
\end{itemize}







\section*{Άσκηση 2  \hfill \small{20 μονάδες}}
Να επιλέξετε τη σωστή απάντηση σε κάθε μία από τις παρακάτω περιπτώσεις:
\begin{enumerate}[1)]
 \item το άθροισμα $13+0$ ισούται με %1
\begin{multicols}{4}
\begin{enumerate}[i)]
 \item $13$
 \item $0$
 \item $14$
 \item $12$
\end{enumerate}
\end{multicols}
 \item το  άθροισμα $91+7$ ισούται με %t2
\begin{multicols}{4}
 \begin{enumerate}[i)]
 \item $-99$
 \item $90+6$
 \item $7+91$
 \item $7+90$
 \end{enumerate}
\end{multicols}
 %-----------------------------------
  \item το άθροισμα $3+(4+6)$ ισούται με %3
\begin{multicols}{4}
\begin{enumerate}[i)]
 \item $3+4+6+6$
 \item $(3+4)+6$
 \item $4+3+6+4$
 \item $4+3+0+3$
\end{enumerate}
\end{multicols}
 \item το άθροισμα $(5+1)+18$ ισούται με %4
\begin{multicols}{4}
\begin{enumerate}[i)]
 \item $5+1+18+1$
 \item $1+5+18+5$
 \item $5+(1+18)+8$
 \item $5+(1+18)$
\end{enumerate}
\end{multicols}
%------------------------------------
 \item το άθροισμα $3+17+0$ ισούται με %5
\begin{multicols}{4}
\begin{enumerate}[i)]
 \item $17+3$
 \item $3+17+7$
 \item $17+3+3$
 \item $4+3+10$
\end{enumerate}
\end{multicols}
%--------------------------------------------
\end{enumerate}

\section*{Άσκηση 3  \hfill \small{} 20 μονάδες}
Να χαρακτηρίσετε καθεμία από τις παρακάτω προτάσεις με τη λέξη "\textbf{Σωστό}``, αν είναι σωστή ή  
"\textbf{Λάθος}`` , αν είναι λανθασμένη:
\begin{enumerate}[1)]
 \item $3+4=4+3$
 \item $7+11+0=0+11$
 \item $2+3+5=2+3+5+0$
 \item $10+8=(8+2)+8$
 \item $6+4+1=6+6$
\end{enumerate}

\section*{Θεωρία - Αφαίρεση Φυσικών Αριθμών \hfill \small{}}
\begin{itemize}
 \item \textbf{Αφαίρεση:} είναι η πράξη με την οποία, όταν δίνονται δύο αριθμοί $Μ$ (μειωτέος) και $Α$ (αφαιρετέος) 
        βρίσκουμε έναν αριθμό $Δ$ (διαφορά) ο οποίος όταν προστεθεί στον $Α$ δίνει το $Μ$,\\
        δηλαδή $Μ=Α+Δ$ και $Δ=Μ-Α$\\
 \textbf{π.χ} $10=14-4$ εδώ μειωτέος:$14$, αφαιρετέος:$4$ και διαφορά:$10$, ισχύει ότι $14=4+10$
 \item Στους φυσικούς αριθμούς ο αφαιρετέος $Α$ πρέπει να είναι πάντα μικρότερος ή ίσος του μειωτέου $Μ$. Σε αντίθετη 
       περίπτωση η πράξη της αφαίρεσης δεν μπορεί να εκτελεστεί (και λέμε ότι δεν ορίζεται).
\end{itemize}


\section*{Άσκηση 4  \hfill \small{} 20 μονάδες}
Να αντιστοιχίσετε κάθε στοιχείο της αριστερής στήλης με ένα στοιχείο της δεξιάς στήλης
\begin{itemize}
\begin{multicols}{2}
 \item $20-7$
 \item $6-12$
 \item $28-8$
 \item $5-3$
 \item $15=11+4$
 \item $20$
 \item $15-4=11$
 \item $13$
 \item δεν ορίζεται η αφαίρεση
 \item $2$
\end{multicols}
\end{itemize}

 









%\let\thefootnote\relax\footnotetext{Επιμέλεια : Κασωτάκη Ειρήνη (ikasotaki@gmail.com),    Σμαραγδάκης Κώστας (kesmarag@gmail.com)} 


\end{document}
